\section{Conclusion}
\label{sec:conclusion}

In this paper, we have presented a minimalist, training-free approach to dynamic model ensembling: \textbf{Parameter-Free Task-Space Projection (PFSR)}. Rejecting the trend of introducing over-parameterized routing layers and complex offline calibration splits, PFSR extracts raw task centroids directly from the frozen classification heads of pre-trained experts to perform direct, training-free dynamic projection coordinates. To investigate the utility of orthogonalizing the task axes, we analyzed an advanced extension: \textbf{L{\"o}wdin-Orthogonalized Task-Space Projection (OTSP)}, which constructs an orthonormal task basis in closed form using the symmetric inverse square root of the Gram overlap matrix.

Through a rigorous 10-seed simulation study, we validated that PFSR and OTSP stably recover the expert ensembling ceiling. Crucially, we proved mathematically and empirically that under symmetric task correlations, OTSP and PFSR are exactly equivalent, while under asymmetric task correlations, OTSP systematically underperforms PFSR due to the Noise Amplification and Noise Spillover Penalties under active representation noise. Consequently, our work establishes a major minimalist win: raw unorthogonalized projection (PFSR) is not only computationally simpler, but systematically more robust than its orthogonalized counterpart under noise. Additionally, we characterized the role of normalization in dynamic ensembling, reframing the unconstrained LinearRouter baseline's instability under sample-wise vectorized streaming ($B=1$) as a pedagogical illustration of the mathematical necessity of the probability-simplex constraint. We also showed that the flat joint classification accuracy of 74.46\% under disjoint subspaces is a symptom of the Orthogonal Masking Effect, where any positive ensembling weight on the correct expert yields the ceiling prediction, and we demonstrated that routing accuracy is the primary informative metric in this regime. Furthermore, we demonstrated that simple zero-initialization of Softmax routing parameters behaves as an implicit maximum-entropy prior that shields parametric routers from small-sample inductive drift.

\subsection{Real-World Evaluation Gap and Future Outlook}
While our high-fidelity 192-dimensional calibrated representation sandbox was a crucial vehicle to study coordinate-level routing dynamics, noise scales, and small-sample inductive overfitting in isolation, we recognize that our evaluation is restricted to a controlled synthetic environment. To bridge the evaluation gap and establish the practical significance of our findings for modern deep learning, several key directions must be explored:

\begin{enumerate}
    \item \textbf{Scaling to High-Dimensional Representations and Generative Models:} Real-world models employ much larger embedding spaces, such as $D=768$ (ViT-Base), $D=1024$ (CLIP), and $D=4096$ (LLaMA-7B). Because the L{\"o}wdin Symmetric Orthogonalization operates on a $K \times K$ Gram overlap matrix (where $K$ is the number of specialist experts in the registry), its offline computational complexity is $O(K^3 + K D)$. This represents an exceptionally lightweight operation that can be executed in under a millisecond even for registries of thousands of specialists, making OTSP highly scalable. Furthermore, to scale this framework to generative models (such as LLMs) that lack standard classification heads, we propose \textbf{Data-Free Centroid Representation (DFCR)}. Instead of relying on class prototype weights $W_k$ to perform SVD, DFCR extracts virtual task prototypes directly from internal layers. Specifically, for each transformer layer $l$, we compute the first principal right-singular vector $v^{(l)} \in \mathbb{R}^D$ of the MLP down-projection weight matrix $W_{\text{down}}^{(l)} \in \mathbb{R}^{D_{\text{mlp}} \times D}$ (or Query-Key-Value projection matrices $W_q, W_k, W_v$). Because attention queries and keys function as bilinear operators ($Q K^T$) rather than standard linear projections, SVD centroids extracted from these matrices must be scaled or analyzed symmetrically to account for attention scaling dynamics. Specifically, instead of extracting independent singular vectors from $W_q$ and $W_k$, one can extract the principal eigenvectors of the combined bilinear query-key mapping $W_q W_k^T$ or scale the individual singular vectors by the standard attention temperature factor $1 / \sqrt{d_k}$ to preserve the attention logit scale across experts. We then define the virtual task centroid $\bar{v}_k$ as the normalized depth-averaged singular vector across selected deep layers of the network: $\bar{v}_k = \text{Normalize}(\frac{1}{|L|}\sum_{l \in L} v^{(l)})$, representing the coordinate direction of maximum semantic variance. These act as intrinsic, zero-shot semantic anchors representing the expert's specialized task manifold without requiring any training or calibration data, making the dynamic projection paradigm fully applicable for LLM merging practitioners.
    \item \textbf{Anisotropic Feature Noise and Covariance Whitening:} Under real-world settings, representation noise is heavily anisotropic (non-spherical) and exhibits severe "narrow-cone" properties in modern embedding spaces. Because our SNR Equivalence and Noise Amplification Penalty proofs assume spherical noise ($\mathbb{E}[\eta_b \eta_b^T] = \sigma^2 I_D$), uncorrected anisotropic noise will skew the L{\"o}wdin transformation's eigenvectors and eigenvalues. Specifically, because L{\"o}wdin orthogonalization utilizes the symmetric inverse square root of the Gram overlap matrix $S_{ij} = q_i \cdot q_j$, any anisotropy in the underlying space introduces phantom correlations between orthogonal task axes, dragging the eigenvectors of $S$ toward the narrow-cone direction. This skews the resulting orthonormal coordinates, disproportionately amplifying noise along the compressed dimensions of the feature manifold and degrading routing accuracy. To address this theoretical gap, we propose integrating an offline covariance-whitening step (e.g., Mahalanobis whitening) using a small calibration set. The empirical covariance matrix $\hat{\Sigma} = \frac{1}{N_{\text{cal}}} \sum_{i=1}^{N_{\text{cal}}} (z_i - \mu)(z_i - \mu)^T + \epsilon I_D$ is estimated from calibration representations $z_i$, and the whitening matrix $\hat{\Sigma}^{-1/2} = U \Lambda^{-1/2} U^T$ is computed via eigendecomposition. Applying this transformation offline to both the extracted SVD centroids ($\bar{v}_k^{\text{white}} = \hat{\Sigma}^{-1/2} \bar{v}_k$) and online to feature representations ($z_b^{\text{white}} = \hat{\Sigma}^{-1/2} z_b$) restores spherical noise properties, successfully neutralizing directional bias, preserving our SNR guarantees, and protecting OTSP's coordinate stability on anisotropic manifolds. To verify this mitigation path, we executed a 2-expert toy simulation under highly anisotropic noise ($\sigma^2_{\text{noisy}} = 1.5$ along the principal centroid dimension versus $\sigma^2_{\text{clean}} = 0.01$ elsewhere). Without whitening, OTSP's routing accuracy falls from 100.0\% (clean) to 77.10\% due to coordinate distortion. However, applying origin-centered second-moment whitening (using the empirical second-moment matrix $R = \mathbb{E}[z z^T]$ to preserve absolute coordinates under the sign-invariant projection) successfully spherizes the noise cloud, restoring OTSP's routing accuracy to 89.45\% (a +12.35\% absolute recovery) and empirically validating the whitening formulation. We explicitly note that while covariance whitening successfully restores coordinate stability, estimating the empirical second-moment or covariance matrix reintroduces a dependency on a small calibration set. This presents a minor trade-off between absolute data-free deployment and anisotropic noise robustness. To minimize this overhead, we emphasize that the covariance matrix is highly general and depends solely on the base model's feature distribution, meaning it can be computed offline once on any generic unlabelled dataset (e.g., a tiny slice of the base model's pre-training corpus) completely independent of the expert tasks. This preserves the task-data-free and label-free nature of our dynamic merging registry, completely avoiding the need for task-specific data splits.
    \item \textbf{Evaluation on Large-Scale Multi-Task Benchmarks:} Future work will evaluate OTSP on standard large-scale multi-task benchmarks like GLUE and MMLU for natural language processing and VTAB for computer vision. This involves extracting SVD centroids from fine-tuned transformer specialists or LoRA adapters (such as BERT, ViT, or LLaMA models fine-tuned on individual tasks) and performing parameter-free dynamic ensembling. Scaling our framework to massive transformer registries containing dozens of LLM adapters represents an exciting and highly promising next step, with the proposed Data-Free Centroid Representation (DFCR) serving as the primary theoretical and mathematical pathway to extract virtual task centroids directly from attention or MLP layers without requiring explicit classification head weights.
    \item \textbf{OOD Safety Gates and Unknown Task Routing:} The current Softmax gating mechanism assumes all test samples belong to one of the $K$ tasks. For Out-of-Distribution (OOD) or unknown task samples, this can lead to incorrect, confident routing to specialized experts. To address this risk, we plan to design OOD safety gates—such as a null-routing state or an entropy threshold on the raw projection coordinates. While a global entropy threshold offers maximum simplicity, unbalanced centroid projection scales across specialists can distort raw coordinates and trigger false rejections. To mitigate this complexity, we propose task-specific adaptive thresholds, where each expert's individual projection score is normalized by its offline self-projection magnitude ($||\bar{v}_k \cdot \bar{v}_k||_2$), or fitting lightweight 1D Gaussian Mixture Models on the projection coordinate distributions over a small validation split. Regarding threshold scaling, when experts exhibit heterogeneous class cardinalities, the entropy of their individual task projection scores scales logarithmically with the number of specialized classes ($O(\log C_k)$). Consequently, the null-routing entropy threshold $\theta_{k}$ for expert $k$ must be dynamically scaled by $\log C_k$ (e.g., setting $\theta_k = \theta_0 \cdot \log C_k$ where $\theta_0$ is a base threshold) to maintain consistent OOD sensitivity across specialists with unbalanced vocabulary scales.
    \item \textbf{Robustness to Class Prototype Variance and Cardinality Imbalance:} In expert registries with heterogeneous class cardinalities (e.g., Expert $A$ classifying $2$ classes for sentiment analysis versus Expert $B$ classifying $1000$ classes for ImageNet), the variance and scale of their prototype distributions will be wildly different. Consequently, the top singular value $\sigma_1^{(k)}$ extracted from the SVD of $W_k$ will scale with the class cardinality $C_k$, introducing systematic routing bias toward experts with larger class vocabularies. To prevent this, we propose three elegant coordinate normalization mechanisms: (1) \textit{Singular Value Rescaling}, where raw coordinates are normalized by their corresponding top singular value, $u'_{k,b} = u_{k,b} / \sigma_1^{(k)}$, scaling the coordinate space by the total energy of the principal direction; (2) \textit{Self-Projection Calibration}, where coordinates are normalized by the expected self-projection magnitude of the specialist's own calibration representations, $u'_{k,b} = u_{k,b} / \mathbb{E}_{z \in \mathcal{D}_k}[|\bar{v}_k \cdot \tilde{z}|]$, calibrating coordinates to represent relative deviations from the specialist's centroid; and (3) \textit{Z-Score Coordinate Standardization}, where raw coordinate distributions are pre-computed over a tiny, task-independent unlabelled dataset and standardized to zero-mean and unit-variance, $u'_{k,b} = (u_{k,b} - \mu_k) / \sigma_k$. These non-parametric scaling solutions put all specialists on an identical coordinate distribution, ensuring perfectly unbiased task-space routing.
\end{enumerate}

\subsection{Concluding Remarks}
Our work serves as a strong reminder of the power of Occam's razor. Simple, closed-form linear algebra can match or exceed the performance of over-engineered, parameter-heavy ensembling pipelines. We call on the machine learning community to relentlessly prune architectural and training complexity, prioritizing elegant mathematical foundations before resorting to optimization-heavy alternatives.
